\documentclass[11pt]{article}
\usepackage[T1]{fontenc}
\usepackage{amsmath,amssymb,amsthm}
\usepackage[margin=1in]{geometry}
\usepackage[hidelinks]{hyperref}

\newtheorem{theorem}{Theorem}
\newtheorem{lemma}[theorem]{Lemma}

\title{The $(k-1)$-Tournament Sch\"utte Boundary:\\
       $f(k-1,k)=k+3$}
\author{Xiaojian Zeng\\
PBC School of Finance, Tsinghua University, 100190\\
\texttt{xiaojianzeng@foxmail.com}}
\date{}

\begin{document}
\maketitle

\begin{abstract}
For integers $m,k\geq 1$, let $f(m,k)$ be the least order of a common
vertex set carrying $m$ tournaments such that every $k$-set is dominated
by one outside vertex in one of the tournaments.  We prove that
\[
    f(k-1,k)=k+3 \qquad (k\geq 3).
\]
The lower bound comes from an auxiliary graph attached to a tournament:
the graph is contained in a star or is exactly a triangle.  This reduces
the only two smaller orders to a source count and an elementary edge-cover
obstruction.  We also give a direct block construction on $k+3$ vertices,
including the endpoint $k=3$.
\end{abstract}

\section{Context and statement}

A \emph{tournament} is a loopless orientation of a complete graph.  Write
$x\to_T y$ when the edge $xy$ is oriented from $x$ to $y$ in a tournament
$T$.  The classical one-tournament domination problem goes back to
Erd\H{o}s~\cite{Erdos1963}; its lower bound was subsequently sharpened by
Szekeres and Szekeres~\cite{SzekeresSzekeres1965}.  Intersecting-family
reformulations have also been useful for a bipartite version of that
problem~\cite{BorowieckiEtAl2003}.

Jeffries~\cite{Jeffries2026} introduced the following multi-tournament
parameter.  For integers $m,k\geq1$, let $f(m,k)$ be the least integer
$n\geq k+1$ for which there are tournaments $T_1,\ldots,T_m$ on one
$n$-element set $V$ satisfying
\begin{equation}\label{eq:property}
  \forall U\in\binom{V}{k}\quad
  \exists i\in\{1,\ldots,m\}\quad
  \exists x\in V\setminus U\quad
  \forall u\in U,\quad x\to_{T_i}u.
\end{equation}
Thus the same outside vertex, in the same tournament, must beat every
member of $U$.  Proposition~2.3 and Theorem~2.2 of
Jeffries~\cite{Jeffries2026} give the upper bound in the theorem below,
and the table there reports equality on this parameter slice for
$k=3,4,5,6$ using a SAT formulation.  The source notes that general
nontrivial lower bounds for $f(m,k)$ are absent, but it does not present
this slice as a numbered open problem; no claim of publication priority is
made here.

Our proof is symbolic and elementary.  Its lower-bound reduction asks
whether a collection of star-contained graphs and triangles can cover the
edges of a complete graph.  This is distinct from the triangle
covering/packing problem studied, for example, in the split-graph setting
by Zeng~\cite{Zeng2026Tuza}: here the triangles are available subgraphs
whose \emph{union} must cover edges.

\begin{theorem}\label{thm:main}
For every integer $k\geq3$,
\[
    f(k-1,k)=k+3.
\]
\end{theorem}

\section{The auxiliary graph}

For a tournament $T$ on a set $V$ of size $n\geq3$, define a simple graph
$H_T$ on $V$ by
\begin{equation}\label{eq:auxiliary}
 ab\in E(H_T)
 \quad\Longleftrightarrow\quad
 a\to_T V\setminus\{a,b\}
 \ \text{or}\ 
 b\to_T V\setminus\{a,b\},
\end{equation}
where $x\to_T W$ means $x\to_T w$ for every $w\in W$.

\begin{lemma}\label{lem:structure}
For every tournament $T$ of order at least three, $H_T$ is contained in a
star or $H_T$ is exactly a triangle.
\end{lemma}

\begin{proof}
First suppose that $T$ has a source $s$, meaning that $s$ beats every
other vertex.  If $ab\in E(H_T)$ and $s\notin\{a,b\}$, then a witnessing
endpoint of $ab$ beats $s$, contradicting that $s$ is a source.  Hence
every edge of $H_T$ is incident with $s$.

Now suppose that $T$ has no source.  If $ab\in E(H_T)$ and $a$ is a
witnessing endpoint in~\eqref{eq:auxiliary}, every possible in-neighbor of
$a$ lies in $\{b\}$.  Since $a$ is not a source, $b\to_T a$; in
particular, the mate of a witnessing endpoint is its unique in-neighbor.
There cannot be two disjoint edges of $H_T$: witnessing endpoints $x$ and
$y$ for two such edges would satisfy both $x\to_T y$ and $y\to_T x$.
Thus the edges of $H_T$ are pairwise intersecting.

It remains only the elementary two-uniform classification.  Empty and
one-edge families are star-contained.  Otherwise choose two edges $ab$
and $ac$.  If not every edge contains $a$, an edge missing $a$ must meet
both chosen edges and hence is $bc$.  Any edge meeting all three of
$ab,ac,bc$ is one of those three.  The family is therefore either
star-contained or exactly the triangle on $\{a,b,c\}$.
\end{proof}

\section{Proof of the boundary}

\begin{proof}[Proof of Theorem~\ref{thm:main}]
Fix $k\geq3$ and put $m=k-1$.

We first rule out order $n=k+1$.  For every $x\in V$, the set
$U_x=V\setminus\{x\}$ has $x$ as its only possible outside witness.
Consequently, covering $U_x$ requires $x$ to be a source in one of the
$m$ tournaments.  A tournament has at most one source, so $m=k-1$
tournaments can account for at most $k-1$ of the $k+1$ choices of $x$.
Thus order $k+1$ is impossible.

Next let $n=k+2$.  The $k$-sets are precisely the complements of pairs,
so~\eqref{eq:property} is equivalent to
\begin{equation}\label{eq:cover}
  E(H_{T_1})\cup\cdots\cup E(H_{T_m})=E(K_n),
  \qquad m=k-1=n-3.
\end{equation}
Use Lemma~\ref{lem:structure} to call an index \emph{star-type} when its
auxiliary graph is contained in a star, and \emph{triangle-type}
otherwise.  In the latter case the auxiliary graph is exactly a triangle.

Suppose there are $a$ star-type indices.  Choose one star center for each
of them (arbitrarily when the auxiliary graph is empty), let $C$ be the
set of distinct chosen centers, and write
\[
    r=|C|,\qquad q=n-r.
\]
Because $r\leq a\leq n-3$, we have $q\geq3$.  After deleting $C$, none of
the star-type graphs supplies an edge.  Hence~\eqref{eq:cover} would force
the remaining $n-3-a$ triangle graphs to cover every edge of $K_q$.  But
\[
 n-3-a\leq n-3-r=q-3,
\]
so their union has at most $3(q-3)$ edges, whereas
\begin{align*}
  2\left(\binom q2-3(q-3)\right)
    &=q^2-7q+18\\
    &=\left(q-\frac72\right)^2+\frac{23}{4}>0.
\end{align*}
This contradiction includes repeated star centers and the endpoint
$q=3$.  Order $k+2$ is therefore impossible.  Since the definition
requires $n\geq k+1$, these are all possible smaller orders, and
$f(k-1,k)\geq k+3$.

For the reverse inequality, partition a set $V$ of size $k+3$ into two
three-element blocks and $k-3$ singleton blocks.  There are exactly
$k-1$ blocks.  For each block $D$, form a tournament $T_D$ by orienting
every edge between $D$ and $V\setminus D$ outward from $D$; if $D$ is a
triple, orient it as a directed $3$-cycle; orient all remaining edges
arbitrarily.  These prescriptions concern disjoint classes of unordered
pairs, so they extend to a tournament.

Let $U\in\binom{V}{k}$.  Either $U$ misses a singleton block, or it
meets one of the triple blocks in at most one vertex.  Indeed, otherwise
$U$ would contain all $k-3$ singleton vertices and at least two vertices
from each triple, giving $|U|\geq k+1$.  If $U$ misses the singleton
$\{x\}$, then $x$ dominates $U$ in $T_{\{x\}}$.  If a triple $D$ is
disjoint from $U$, any vertex of $D$ dominates $U$ in $T_D$.  Finally, if
$D\cap U=\{u\}$, the directed cycle on $D$ contains a vertex
$x\in D\setminus U$ with $x\to_{T_D}u$, and the outward cross-block
orientation makes $x$ beat every vertex of $U\setminus D$.  Thus
property~\eqref{eq:property} holds on $k+3$ vertices.  The same argument
works when $k=3$, when there are no singleton blocks.  Hence
$f(k-1,k)\leq k+3$, completing the proof.
\end{proof}

\section*{Computational and AI assistance}

OpenAI's Codex system assisted with literature searching, adversarial proof
checking, preparation of a finite diagnostic script, and language editing.
The script exhaustively checked Lemma~\ref{lem:structure} for labeled
tournaments of orders $3$ through $6$, the block construction for
$3\leq k\leq12$, and the displayed inequality for $3\leq q\leq100$.
These finite checks are sanity checks only; no computational output is used
to prove the quantified theorem.  The author is responsible for the
statement, proof, citations, and manuscript.

\bibliographystyle{plain}
\bibliography{references}

\end{document}
